% !TEX TS-program = pdflatex
% academicthesis-user-guide.tex
% User guide of the academicthesis class for the LUISS and Politecnico di
% Milano packages. Compile from this folder with pdfLaTeX (twice, for the
% table of contents and cross-references):
%   pdflatex academicthesis-user-guide.tex
% It is a standalone document: it does not need the academicthesis class.
\documentclass[11pt,a4paper]{article}
\newcommand{\doctitleshort}{academicthesis --- User guide}
% academicthesis-docstyle.tex
% Shared preamble of the academicthesis documentation (introduction and
% user guide). \input it right after \documentclass. Standalone: the
% documentation does not need the academicthesis class itself.
\usepackage[T1]{fontenc}
\usepackage[utf8]{inputenc}
\usepackage{lmodern}
\usepackage[english]{babel}
\usepackage[a4paper, margin=2.6cm, headheight=14pt]{geometry}
\usepackage{microtype}
\usepackage{xcolor}
\usepackage{graphicx}
\usepackage{array}
\usepackage{tabularx}
\usepackage{booktabs}
\usepackage{longtable}
\usepackage{listings}
\usepackage{enumitem}
\usepackage{titlesec}
\usepackage{fancyhdr}
\usepackage{tcolorbox}
\tcbuselibrary{skins, breakable}
\usepackage{tikz}
\usetikzlibrary{positioning, arrows.meta, fit}
\usepackage{url}
\usepackage[colorlinks=true, linkcolor=docaccent, urlcolor=docaccent,
            citecolor=docaccent]{hyperref}

% ---- Colours ---------------------------------------------------------
\definecolor{docaccent}{RGB}{31,73,125}   % neutral blue for headings
\definecolor{luissblue}{RGB}{21,57,108}   % BlueLuiss, as in the package
\definecolor{bluepoli}{cmyk}{0.4,0.1,0,0.4} % bluepoli, as in the package
\definecolor{docgrey}{gray}{0.45}

% ---- Page and paragraphs --------------------------------------------
\setlength{\parindent}{0pt}
\setlength{\parskip}{0.55\baselineskip}
\setlength{\emergencystretch}{3em}
\setlist{nosep, topsep=3pt, leftmargin=1.5em}
\renewcommand{\arraystretch}{1.25}

\titleformat{\section}{\color{docaccent}\normalfont\Large\bfseries}{\thesection}{0.8em}{}
\titleformat{\subsection}{\color{docaccent}\normalfont\large\bfseries}{\thesubsection}{0.8em}{}
\titleformat{\subsubsection}{\normalfont\normalsize\bfseries}{\thesubsubsection}{0.8em}{}
\titlespacing*{\section}{0pt}{2.2ex plus 1ex}{1ex}

\pagestyle{fancy}
\fancyhf{}
\fancyhead[L]{\small\color{docgrey}\doctitleshort}
\fancyhead[R]{\small\color{docgrey}\nouppercase{\leftmark}}
\fancyfoot[C]{\small\thepage}
\renewcommand{\headrulewidth}{0.3pt}
\renewcommand{\sectionmark}[1]{\markboth{\thesection\quad #1}{}}
\fancypagestyle{plain}{\fancyhf{}\renewcommand{\headrulewidth}{0pt}\fancyfoot[C]{\small\thepage}}
\providecommand{\doctitleshort}{academicthesis}

% ---- Names in running text and tables -------------------------------
\newcommand{\cs}[1]{\texttt{\textbackslash#1}}   % a command
\newcommand{\opt}[1]{\texttt{#1}}                 % an option / key / value
\newcommand{\env}[1]{\texttt{#1}}                 % an environment
\newcommand{\pkg}[1]{\textsf{#1}}                 % a package or class
\DeclareUrlCommand{\file}{\urlstyle{tt}}          % a file name (breaks at - . /)
\renewcommand{\UrlBreaks}{\do\-\do\.\do\/}
\newcolumntype{Y}{>{\raggedright\arraybackslash}X}
\newcolumntype{L}[1]{>{\raggedright\arraybackslash}p{#1}}
\newcommand{\req}{\textbf{required}}
\newcommand{\optional}{optional}
\newcommand{\na}{\textcolor{docgrey}{---}}

% ---- Code listings ----------------------------------------------------
\lstset{
  basicstyle=\ttfamily\small,
  keywordstyle=\color{docaccent}\bfseries,
  commentstyle=\color{docgrey}\itshape,
  language=[LaTeX]TeX,
  morekeywords={thesissetup, frontpage, luissthesissetup, polimithesissetup,
    major, dedication, frontmatter, mainmatter, backmatter, addbibresource,
    printbibliography, blankpage},
  frame=single, rulecolor=\color{docaccent!35},
  backgroundcolor=\color{docaccent!4},
  columns=fullflexible, keepspaces=true, upquote=true,
  xleftmargin=4pt, xrightmargin=4pt, aboveskip=8pt, belowskip=8pt,
  literate={à}{{\`a}}1 {è}{{\`e}}1 {é}{{\'e}}1 {ò}{{\`o}}1,
}
\lstdefinestyle{shell}{language={}, keywordstyle={}, commentstyle={}}

% ---- Boxes -----------------------------------------------------------
% \begin{notebox}{Title} ... \end{notebox}
\newtcolorbox{notebox}[2][docaccent]{enhanced, breakable,
  colback=#1!5, colframe=#1!70, boxrule=0.6pt, arc=2pt,
  left=8pt, right=8pt, top=6pt, bottom=6pt,
  fonttitle=\bfseries\small, coltitle=white, title={#2}}

% A framed page preview: \preview[<width>]{<file>}{<caption>}
\newcommand{\preview}[3][0.42\textwidth]{%
  \begin{minipage}[t]{#1}\centering
    {\setlength{\fboxsep}{0pt}\setlength{\fboxrule}{0.5pt}%
     \fcolorbox{docgrey!50}{white}{\includegraphics[width=\dimexpr\linewidth-1pt\relax]{#2}}}\par
    \vspace{4pt}{\small #3}
  \end{minipage}}

\graphicspath{{figures/}}


\hypersetup{pdftitle={academicthesis: user guide for the LUISS and Politecnico di Milano packages},
            pdfauthor={Luca Salvatore}}

\begin{document}

% ======================================================================
% Title page
% ======================================================================
\hypersetup{pageanchor=false}
\begin{titlepage}
  \thispagestyle{empty}
  \vspace*{3cm}
  {\color{docaccent}\rule{\linewidth}{1.2pt}}\par\vspace{14pt}
  {\fontsize{30}{36}\selectfont\bfseries\color{docaccent} academicthesis\par}
  \vspace{10pt}
  {\LARGE User guide\par}
  \vspace{6pt}
  {\Large Writing a thesis with the LUISS and Politecnico di Milano packages\par}
  \vspace{14pt}
  {\color{docaccent}\rule{\linewidth}{1.2pt}}\par
  \vspace{1.5cm}
  \begin{minipage}{0.8\linewidth}
    This guide is for students who write their thesis with
    \pkg{academicthesis}. It explains how to compile a thesis, every class
    option and data field, what each university and format requires, and
    gives a complete example for each layout.
  \end{minipage}
  \vfill
  {\large Luca Salvatore\par}
  {\color{docgrey} Class version 0.2.0 \quad\textbullet\quad October 2026\par}
\end{titlepage}
\hypersetup{pageanchor=true}

{\setlength{\parskip}{0pt}\tableofcontents}
\clearpage

% ======================================================================
\section{About this guide}

\pkg{academicthesis} is a \LaTeX{} class that typesets a thesis in the
format prescribed by a university. You choose the university, the degree
and, where a university offers several, the format; you give the data of
your thesis once; the class produces the cover page, front matter and page
layout. This guide covers the two university packages documented so far:

\begin{itemize}
  \item \pkg{academicthesis-luiss}, for LUISS Guido Carli (bachelor's,
    master's and doctoral theses), Section~\ref{sec:luiss};
  \item \pkg{academicthesis-polimi}, for Politecnico di Milano (classical,
    article and Executive Summary formats, and the doctoral dissertation),
    Section~\ref{sec:polimi}.
\end{itemize}

Sections~\ref{sec:install}--\ref{sec:bib} describe what is common to every
university; read them first, then the section of your university.
Section~\ref{sec:errors} lists the most common error messages and how to
fix them, and Appendix~\ref{sec:cheatsheet} is a one-page reference.

\paragraph{Conventions.} Commands are written \cs{thesissetup}, options and
keys \opt{format=classical}, files \file{example.tex}. ``Required'' means
that compilation stops with an error naming the field when it is missing;
``optional'' means it is printed only when given.

% ======================================================================
\section{Installing and compiling}
\label{sec:install}

\subsection{Requirements}

\begin{itemize}
  \item A recent \TeX{} distribution: \TeX{}~Live or Mac\TeX{} 2023 or
    later, or an up-to-date MiK\TeX. The class is written in \LaTeX3 and
    needs a modern kernel.
  \item \textbf{pdf\LaTeX} is the reference engine and the one used by all
    the examples. Xe\LaTeX{} and Lua\LaTeX{} are also accepted (the class then
    loads \pkg{fontspec}).
  \item For the bibliography, \pkg{biblatex} with \pkg{biber} (both part of
    a full \TeX{} distribution). Without them the class falls back to
    \pkg{natbib} and Bib\TeX{} (Section~\ref{sec:bib}).
\end{itemize}

\subsection{Which files you need}

The class is a collection of small packages that load one another by name:

\begin{tabularx}{\textwidth}{@{}l Y@{}}
\toprule
\textbf{Files} & \textbf{Needed for} \\
\midrule
\file{academicthesis.cls}, \file{academicthesis-kernel.sty} & Always (the core). \\
\file{academicthesis-english.ldf}, \file{-italian.ldf}, \file{-french.ldf} & The languages you use. \\
\file{universities/luiss/*} & A LUISS thesis. \\
\file{universities/polimi/*} & A Politecnico di Milano thesis. \\
\bottomrule
\end{tabularx}

Any of these three set-ups works:

\begin{enumerate}
  \item \textbf{Inside the repository} (simplest). Put your thesis in the
    \file{showcases/} folder, or copy a showcase and keep its first line,
    \verb|\input{academicthesis-repo-paths}|, which tells \LaTeX{} where the
    package folders are. Then compile with any editor, or with
    \begin{lstlisting}[style=shell]
cd showcases
latexmk my-thesis.tex
    \end{lstlisting}
  \item \textbf{Next to your thesis.} Copy the files listed above into the
    folder of your thesis (no sub-folders needed), and remove the
    \verb|\input{academicthesis-repo-paths}| line. This is also the way to
    use the class on Overleaf: upload the files to the project.
  \item \textbf{Installed in your personal \TeX{} tree}, for example under
    \file{~/Library/texmf/tex/latex/academicthesis/} (macOS) or
    \file{~/texmf/tex/latex/academicthesis/} (Linux). Every thesis on the
    computer can then use the class with no extra line.
\end{enumerate}

\begin{notebox}{Editors}
Texifier/Texpad, TeXShop, TeXstudio, VS Code (LaTeX Workshop) and Overleaf
all work. In Texifier/Texpad, set the typesetter to \emph{TeX Live
(pdfLaTeX)}: its built-in engine is too old for this class.
\end{notebox}

\subsection{Compiling}

Run pdf\LaTeX{} twice (for the table of contents), and \pkg{biber} between
the two runs when you have a bibliography. \texttt{latexmk} does all of this
by itself:

\begin{lstlisting}[style=shell]
latexmk -pdf my-thesis.tex
\end{lstlisting}

% ======================================================================
\section{Anatomy of a thesis}
\label{sec:anatomy}

Every thesis has the same four parts, whatever the university:

\begin{lstlisting}
% 1. Class options: university, degree, format, language
\documentclass[university=luiss, master, language=english]{academicthesis}

% 2. Your own packages (optional)
\usepackage{amssymb}
\usepackage{hyperref}

% 3. The data of the thesis
\thesissetup{
  title        = {Title of the Thesis},
  candidate    = {Name Surname},
  degree       = {Economics and Management},
  chair        = {Corporate Finance},
  supervisor   = {Prof. Name Surname},
  academicyear = {2025/2026},
}
\addbibresource{references.bib}

% 4. The document
\begin{document}
\frontpage                      % cover page (and second page, if any)

\frontmatter                    % roman page numbers
\begin{abstractthesis}[keywords = {first, second}]
  The abstract.
\end{abstractthesis}
\tableofcontents

\mainmatter                     % arabic page numbers
\chapter{Introduction}
...
\printbibliography

\backmatter
\begin{acknowledgements}
  Thanks to ...
\end{acknowledgements}
\end{document}
\end{lstlisting}

\cs{frontpage} prints the cover page of your university, degree and
format from the data given in \cs{thesissetup}. Everything after it is
ordinary \LaTeX: chapters, sections, figures, tables and so on.

% ======================================================================
\section{Class options}
\label{sec:options}

All options go between the square brackets of
\cs{documentclass}\opt{[...]\{academicthesis\}}.

\begin{tabularx}{\textwidth}{@{}L{3.3cm} L{3.4cm} Y@{}}
\toprule
\textbf{Option} & \textbf{Values} & \textbf{Meaning} \\
\midrule
\opt{university} (or \opt{uni}) & \opt{luiss}, \opt{polimi} & The university whose format is used. Required. \\
degree & \opt{bachelor}, \opt{master}, \opt{phd} & The degree. Written on its own, without \opt{degree=}. Default: \opt{bachelor}. \\
\opt{format} & depends on the university & The layout, for universities that offer several (PoliMi: \opt{classical}, \opt{article}, \opt{summary}, \opt{phd}). LUISS has a single layout and does not accept \opt{format}. \\
\opt{language} (or \opt{lang}) & \opt{english}, \opt{italian}, \opt{french} & Main language of the thesis: every printed label, and hyphenation. Default: \opt{english}. Also accepted: \opt{lang=en}, \opt{lang=it}, \opt{lang=fr}, or just \opt{italian}. \\
\opt{other-languages} & e.g.\ \opt{\{italian,french\}} & Other languages used by single blocks (Section~\ref{sec:blocks}), so that their hyphenation is right. PoliMi always loads English and Italian. \\
\opt{secondpage}, \opt{no-secondpage} & & Print or omit the second page after the cover (Section~\ref{sec:secondpage}). Default: on for LUISS, off for PoliMi. \\
\opt{discussion}, \opt{no-discussion} & & Whether the thesis has been defended; it changes the second page and makes \opt{examdate} required. Default: off. \\
\opt{biblatex-options} & e.g.\ \opt{\{style=apa\}} & Options passed to \pkg{biblatex} (Section~\ref{sec:bib}). \\
anything else & e.g.\ \opt{11pt}, \opt{oneside}, \opt{twoside} & Passed to the base class (\pkg{book} or \pkg{article}). Your choice wins over the university's default. \\
\bottomrule
\end{tabularx}

% ======================================================================
\section{The data of your thesis}
\label{sec:data}

All the data printed on the cover and on the second page is given with one
command in the preamble, \cs{thesissetup}, as a list of
\opt{key = \{value\}} pairs. The key names are the same at every university;
only the printed labels change (for example LUISS prints ``Supervisor'',
PoliMi ``Advisor'').

\begin{longtable}{@{}L{3.2cm} L{11.2cm}@{}}
\toprule
\textbf{Key} & \textbf{What to write} \\
\midrule
\endhead
\opt{title} & Title of the thesis. \\
\opt{subtitle} & Subtitle, printed under the title when given. \\
\opt{candidate} & Your name, as it should be printed: \opt{\{Name Surname\}}. \\
\opt{idnumber} & Student ID number (\emph{matricola}). \\
\opt{gender} & \opt{male} or \opt{female} (also \opt{m}, \opt{f}). Only changes Italian labels: ``Candidato'' / ``Candidata''. \\
\opt{degree} & The degree programme, without ``Course of'' or ``Master's Thesis in'', which the cover adds: \opt{\{Economics and Management\}}. \\
\opt{chair} & LUISS: the chair (\emph{cattedra}) of the thesis, without ``Chair of''. \\
\opt{department} & The department: LUISS PhD without ``Department of'' (the cover adds it); PoliMi PhD in full, as it should be printed. \\
\opt{phdtitle} & The doctoral programme: LUISS without ``PhD in''; PoliMi in full (\opt{\{Doctoral Program in \dots\}}). \\
\opt{cycle} & PhD cycle, e.g.\ \opt{XXXVIII}. \\
\opt{academicyear} & Academic year, written as your university writes it: \opt{2025/2026}, \opt{2025-26}. \\
\opt{supervisor} & One supervisor (advisor), with title: \opt{\{Prof. Name Surname\}}. \\
\opt{supervisors} & Several supervisors, separated by commas: \opt{\{Prof. A, Prof. B\}}. \\
\opt{cosupervisor}, \opt{cosupervisors} & Co-supervisor(s) (co-advisors), in the same way. \\
\opt{tutor}, \opt{tutors} & PoliMi PhD: tutor(s). \\
\opt{coordinator} & PoliMi PhD: chair (coordinator) of the doctoral programme. \\
\opt{abstract}, \opt{keywords} & PoliMi \opt{article} format: the English abstract and its keywords, printed on page~1. \\
\opt{copyyear} & Second page: year of the copyright notice. \\
\opt{authoremail} & Second page: contact e-mail. \\
\opt{website} & Second page: website. \\
\opt{examdate} & Second page, with the \opt{discussion} option: date of the defence. \\
\opt{examiner}, \opt{examiners} & Second page, with \opt{discussion}: members of the committee, comma-separated. \\
\bottomrule
\end{longtable}

\begin{notebox}{One person or several}
\opt{supervisor} takes one name and keeps any comma in it;
\opt{supervisors} splits its value at the commas. The label adapts by
itself: ``Supervisor'' for one name, ``Supervisors'' for two or more. The
same holds for \opt{cosupervisor(s)}, \opt{tutor(s)} and
\opt{examiner(s)}. Write academic titles (``Prof.'', ``Dr.'') yourself.
\end{notebox}

\cs{thesissetup} can be used several times; a later value replaces an
earlier one. Which keys are required depends on the university, the degree
and the format: see the tables in Sections~\ref{sec:luiss}
and~\ref{sec:polimi}. Keys that a layout does not use are simply ignored.

% ======================================================================
\section{Front-matter blocks}
\label{sec:blocks}

The core class provides four blocks. Each university gives them its own
look, but they are written the same way everywhere.

\begin{tabularx}{\textwidth}{@{}l l Y@{}}
\toprule
\textbf{Block} & \textbf{Options} & \textbf{Default heading (English)} \\
\midrule
\cs{begin}\opt{\{abstractthesis\}} & \opt{language}, \opt{keywords}, \opt{title} & Abstract \\
\cs{begin}\opt{\{originality\}} & \opt{language}, \opt{title} & Statement of Originality \\
\cs{begin}\opt{\{acknowledgements\}} & \opt{language}, \opt{title} & Acknowledgments \\
\cs{dedication}\opt{[...]\{text\}} & \opt{language} & (none: the text on its own page) \\
\bottomrule
\end{tabularx}

\begin{itemize}
  \item \opt{language=} writes the block in another language: its heading,
    the ``Keywords'' label and the hyphenation follow that language.
    Without it, the block uses the main language of the thesis. List the
    extra language in the \opt{other-languages} class option (not needed
    for Italian and English at PoliMi).
  \item \opt{keywords=} prints the keywords after the abstract.
  \item \opt{title=} replaces the default heading.
\end{itemize}

Abstracts, statements of originality and acknowledgements appear in the
table of contents. For example, an English thesis with an Italian abstract:

\begin{lstlisting}
\documentclass[university=luiss, master, language=english,
               other-languages={italian}]{academicthesis}
...
\begin{abstractthesis}[keywords = {corporate finance, valuation}]
  The English abstract.
\end{abstractthesis}
\begin{abstractthesis}[language = italian,
                       keywords = {finanza aziendale, valutazione}]
  L'abstract in italiano.
\end{abstractthesis}
\dedication{To my family.}
\end{lstlisting}

% ======================================================================
\section{The second page}
\label{sec:secondpage}

Some universities print, right after the cover, a page with the examination
details and a copyright notice. \cs{frontpage} adds it when it is on: by
default for LUISS, not for PoliMi. Switch it with the class options
\opt{secondpage} and \opt{no-secondpage}.

\begin{itemize}
  \item \textbf{Top of the page.} Without the \opt{discussion} option:
    ``The present thesis has not been defended.'' With \opt{discussion}:
    the date of the defence (\opt{examdate}, then required) and the
    committee (\opt{examiners}).
  \item \textbf{Bottom of the page.} The title, ``\textcopyright{}
    \opt{copyyear} \opt{candidate}'', a reservation of rights, the author,
    and the contact e-mail and website when given (\opt{authoremail},
    \opt{website}).
\end{itemize}

Give \opt{copyyear} whenever the second page is printed, or the notice
will have no year.

% ======================================================================
\section{Bibliography}
\label{sec:bib}

The class loads \pkg{biblatex} with \pkg{biber} for you. The style is, in
order of priority: the one you give with the class option
\opt{biblatex-options}; the university's default (LUISS: author--year);
otherwise numeric.

\begin{lstlisting}
\documentclass[university=polimi, master,
               biblatex-options={style=numeric-comp, sorting=none}]{academicthesis}
\addbibresource{references.bib}   % in the preamble
...
\printbibliography                % where the bibliography goes
\end{lstlisting}

Do not load \pkg{biblatex} again yourself. If \pkg{biblatex} is not
installed, the class warns and uses \pkg{natbib} with Bib\TeX{}, keeping
the same two commands, so the thesis still compiles.

% ======================================================================
\section{LUISS Guido Carli}
\label{sec:luiss}

\begin{center}
  \preview[0.28\textwidth]{luiss-master-cover}{Bachelor's / master's cover}\hfill
  \preview[0.28\textwidth]{luiss-secondpage}{Second page (with \opt{discussion})}\hfill
  \preview[0.28\textwidth]{luiss-phd-cover}{PhD cover}
\end{center}

\subsection{At a glance}

\begin{tabularx}{\textwidth}{@{}l Y@{}}
\toprule
Class line & \cs{documentclass}\opt{[university=luiss, master]\{academicthesis\}} \\
Degrees & \opt{bachelor}, \opt{master}, \opt{phd} \\
Formats & one; the \opt{format} option is not used \\
Base class & \pkg{book}, with the class defaults (add \opt{twoside}, \opt{12pt}, \dots{} as you wish) \\
Second page & on (\opt{no-secondpage} to omit it) \\
Bibliography & author--year (\pkg{biblatex}, \pkg{biber}) \\
Languages & English, Italian or French labels, after the \opt{language} option \\
\bottomrule
\end{tabularx}

\subsection{What the cover shows}

\textbf{Bachelor's and master's thesis}, from top to bottom: the LUISS
logo; ``Course of \opt{degree}'' (Italian: ``Corso di Laurea''); the major,
when given; ``Chair of \opt{chair}'' (``Cattedra''); the title in LUISS blue
and the subtitle; the supervisor(s) on the left, the co-supervisor(s) on the
right and the candidate with the student ID below them (or, without a
co-supervisor, supervisor on the left and candidate on the right); the
academic year at the bottom.

\textbf{PhD thesis}: the same page, with ``Department of \opt{department}''
under the logo, ``PhD in \opt{phdtitle}'' in place of the course, and
``Cycle \opt{cycle}'' in place of the chair.

\subsection{Thesis data}

\begin{tabularx}{\textwidth}{@{}L{3.2cm} Y c c@{}}
\toprule
\textbf{Key} & \textbf{Printed as} & \textbf{Bachelor / master} & \textbf{PhD} \\
\midrule
\opt{title} & Title, in LUISS blue & \req & \req \\
\opt{subtitle} & Under the title & \optional & \optional \\
\opt{candidate} & ``Candidate'' column & \req & \req \\
\opt{idnumber} & Under the candidate's name & \optional & \optional \\
\opt{degree} & ``Course of \dots'' & \req & \na \\
\opt{chair} & ``Chair of \dots'' & \req & \na \\
\opt{department} & ``Department of \dots'' & \na & \req \\
\opt{phdtitle} & ``PhD in \dots'' & \na & \req \\
\opt{cycle} & ``Cycle \dots'' & \na & \req \\
\opt{supervisor(s)} & ``Supervisor(s)'' column & \req & \req \\
\opt{cosupervisor(s)} & ``Co-Supervisor(s)'' column & \optional & \optional \\
\opt{academicyear} & ``Academic Year'' & \req & \req \\
\opt{gender} & ``Candidato'' / ``Candidata'' (Italian) & \optional & \optional \\
\opt{copyyear}, \opt{authoremail}, \opt{website} & Second page & recommended & recommended \\
\opt{examdate}, \opt{examiners} & Second page & with \opt{discussion} & with \opt{discussion} \\
\bottomrule
\end{tabularx}

\textbf{The major} (bachelor's and master's) is LUISS-specific and is set
with its own command, or with \cs{luissthesissetup}; either line below
works. It is printed as written, between the course and the chair.
\begin{lstlisting}
\major{Major in Corporate Finance}
\luissthesissetup{major = {Major in Corporate Finance}}
\end{lstlisting}

\subsection{Adjusting the cover: \texorpdfstring{\cs{luissthesissetup}}{\\luissthesissetup}}

\begin{tabularx}{\textwidth}{@{}l l Y@{}}
\toprule
\textbf{Key} & \textbf{Default} & \textbf{Effect} \\
\midrule
\opt{logo-file} & \opt{logoluissblue} & Logo image (file name without extension). \opt{luiss} is also shipped. \\
\opt{logo-width} & \opt{0.55\textbackslash textwidth} & Width of the logo. \\
\opt{logo-align} & \opt{left} & \opt{left}, \opt{center} or \opt{right}. \\
\opt{title-align} & \opt{left} & Title and subtitle. \\
\opt{course-align} & \opt{left} & Course and chair block (PhD: programme and cycle). \\
\opt{major-align} & \opt{left} & The major line. \\
\opt{department-align} & \opt{left} & PhD: the department line. \\
\opt{major} & & The major (see above). \\
\bottomrule
\end{tabularx}

The alignment keys other than \opt{logo-align} also accept the short forms
\opt{l}, \opt{c}, \opt{r}. Example: \verb|\luissthesissetup{logo-align=center, title-align=center}|.

\subsection{Example: bachelor's or master's thesis}

\begin{lstlisting}
\documentclass[university=luiss, master, language=english, twoside,
               discussion, other-languages={italian}]{academicthesis}

\major{Major in Corporate Finance}

\thesissetup{
  title        = {Title of the Thesis},
  subtitle     = {An Optional Subtitle},
  candidate    = {Name Surname},
  idnumber     = {000000},
  degree       = {Economics and Management},
  chair        = {Corporate Finance},
  supervisor   = {Prof. Name Surname},
  cosupervisor = {Prof. Name Surname},
  academicyear = {2025/2026},
  copyyear     = {2026},
  authoremail  = {name.surname@example.com},
  examdate     = {12 July 2026},              % because of `discussion'
  examiners    = {Prof. A, Prof. B, Prof. C},
}
\addbibresource{references.bib}

\begin{document}
\frontpage

\frontmatter
\begin{abstractthesis}[keywords = {first, second}]
  The abstract.
\end{abstractthesis}
\begin{abstractthesis}[language = italian, keywords = {primo, secondo}]
  Il sommario in italiano.
\end{abstractthesis}
\dedication{To my family.}
\tableofcontents

\mainmatter
\chapter{Introduction}
...
\printbibliography

\backmatter
\begin{acknowledgements}
  Thanks to ...
\end{acknowledgements}
\end{document}
\end{lstlisting}

For a bachelor's thesis, write \opt{bachelor} instead of \opt{master}.
Without \opt{discussion}, leave out \opt{examdate} and \opt{examiners}; the
second page then says that the thesis has not been defended.

\subsection{Example: PhD thesis}

\begin{lstlisting}
\documentclass[university=luiss, phd, language=english, twoside]{academicthesis}

\thesissetup{
  title        = {Title of the Dissertation},
  candidate    = {Name Surname},
  idnumber     = {000000},
  department   = {Economics and Finance},   % "Department of ..."
  phdtitle     = {Economics},               % "PhD in ..."
  cycle        = {XXXVIII},                 % "Cycle ..."
  supervisors  = {Prof. Name Surname},
  academicyear = {2025/2026},
  copyyear     = {2026},
}

\begin{document}
\frontpage
\frontmatter
\begin{abstractthesis} The abstract. \end{abstractthesis}
\tableofcontents
\mainmatter
\chapter{Introduction}
...
\end{document}
\end{lstlisting}

Complete versions are in \file{showcases/example-luiss-bachelor.tex} and
\file{showcases/example-luiss-phd.tex}.

% ======================================================================
\section{Politecnico di Milano}
\label{sec:polimi}

\begin{notebox}[bluepoli]{The PoliMi layouts are Politecnico di Milano's design}
The \opt{classical}, \opt{article} and \opt{summary} formats reproduce the
official thesis templates of the School of Industrial and Information
Engineering (Scuola 3I) of Politecnico di Milano, revision of October 2021
(\emph{Thesis classical format}, \emph{Thesis article format},
\emph{Executive Summary}). The \opt{phd} format reproduces the PoliMi
doctoral dissertation template (\file{phdthesis.sty},
\file{phdtitle.sty}). The page designs, colours, logos, seal and
``raggiera'' belong to Politecnico di Milano; this package re-implements
them on the \pkg{academicthesis} core and is not the original author of
the design. Check your school's current rules before submitting.
\end{notebox}

\begin{center}
  \preview[0.225\textwidth]{polimi-classical-cover}{\opt{classical}}\hfill
  \preview[0.225\textwidth]{polimi-article-cover}{\opt{article}}\hfill
  \preview[0.225\textwidth]{polimi-summary-cover}{\opt{summary}}\hfill
  \preview[0.18\textwidth]{polimi-phd-cover}{\opt{phd}}
\end{center}

\subsection{At a glance}

\begin{tabularx}{\textwidth}{@{}l Y l l@{}}
\toprule
\textbf{Format} & \textbf{What it is} & \textbf{Degrees} & \textbf{Base class} \\
\midrule
\opt{classical} & The full thesis as a book. Default for bachelor and master. & bachelor, master & \pkg{book}, 12pt, two-sided, A4 \\
\opt{article} & The thesis as a paper, with the abstract on page~1. & bachelor, master & \pkg{article}, 11pt, A4 \\
\opt{summary} & The Executive Summary, in two columns. & bachelor, master & \pkg{article}, 11pt, 2 col., A4 \\
\opt{phd} & The doctoral dissertation. Default for \opt{phd}. & phd & \pkg{book}, 10pt, 17$\times$24\,cm \\
\bottomrule
\end{tabularx}

\begin{itemize}
  \item Class line: \cs{documentclass}\opt{[university=polimi, master, format=classical]}\linebreak[0]\opt{\{academicthesis\}}.
  \item No second page (add it with \opt{secondpage} if you need it).
  \item English and Italian are always loaded, so the Italian abstract
    needs no \opt{other-languages}.
  \item The package sets no bibliography style; the default is numeric
    (Section~\ref{sec:bib}).
\end{itemize}

\subsection{Thesis data}

\begin{tabularx}{\textwidth}{@{}L{3.1cm} Y c c c@{}}
\toprule
\textbf{Key} & \textbf{Printed as (English / Italian)} & \textbf{classical} & \textbf{article} & \textbf{summary} \\
\midrule
\opt{title} & Title, in blue & \req & \req & \req \\
\opt{subtitle} & Under the title & \optional & \optional & \optional \\
\opt{candidate} & Author / Candidato & \req & \req & \req \\
\opt{idnumber} & Student ID / Matricola & \req & \req & \na \\
\opt{degree} & after ``Master's Thesis in'' / ``Tesi di Laurea Magistrale in'' & \req & \req & \req \\
\opt{supervisor(s)} & Advisor(s) / Relatore, Relatori & \req & \req & \req \\
\opt{cosupervisor(s)} & Co-advisor(s) / Correlatore, Correlatori & \optional & \optional & \optional \\
\opt{academicyear} & Academic Year / Anno Accademico & \req & \req & \req \\
\opt{abstract} & English abstract on page~1 & \na & \req & \na \\
\opt{keywords} & Keywords of that abstract, in a blue box & \na & \optional & \na \\
\opt{gender} & Candidato / Candidata (Italian) & \optional & \optional & \optional \\
\bottomrule
\end{tabularx}

For a bachelor's thesis the cover reads ``Bachelor's Thesis in'' (``Tesi di
Laurea Triennale in'').

\subsubsection*{PhD (\texorpdfstring{\opt{format=phd}}{format=phd})}

\begin{tabularx}{\textwidth}{@{}L{3.1cm} Y c@{}}
\toprule
\textbf{Key} & \textbf{Printed as (English / Italian)} & \textbf{Required} \\
\midrule
\opt{title}, \opt{subtitle} & Title; subtitle when given & title \\
\opt{candidate} & Under ``Doctoral Dissertation of'' / ``Tesi di Dottorato di'' & \req \\
\opt{department} & In full, e.g.\ ``Department of Electronics, Information and Bioengineering'' & \req \\
\opt{phdtitle} & In full, e.g.\ ``Doctoral Program in Information Technology'' & \req \\
\opt{supervisor(s)} & Supervisor(s) / Relatore, Relatori & \req \\
\opt{tutor(s)} & Tutor(s) / Tutor & \optional \\
\opt{coordinator} & The Chair of the Doctoral Program / Il Coordinatore del Dottorato & \optional \\
\opt{academicyear}, \opt{cycle} & Bottom line: ``2025--2026 -- Cycle XXXVIII'' & \req \\
\bottomrule
\end{tabularx}

\subsection{Options: \texorpdfstring{\cs{polimithesissetup}}{\\polimithesissetup}}

\begin{tabularx}{\textwidth}{@{}l l Y@{}}
\toprule
\textbf{Key} & \textbf{Default} & \textbf{Effect} \\
\midrule
\opt{logo-file} & \opt{logopolimi3i} & Logo on the cover or first page (file name without extension). The \opt{phd} format uses the round seal, \opt{logopolimiseal}, instead. \\
\opt{watermark} & \opt{true} & The ``raggiera'' in the background of the cover, of page~1 and of blank pages. \opt{false} removes it. Not used by \opt{phd}. \\
\opt{paper} & \opt{17x24} & \opt{phd} only: \opt{17x24} (17$\times$24\,cm, as in the template) or \opt{a4}. \\
\bottomrule
\end{tabularx}

\subsection{The \texorpdfstring{\opt{classical}}{classical} format}

A4, 12pt, two-sided, 1.5 line spacing, no paragraph indent. The cover has
the Scuola 3I logo, the title in blue, ``Master's Thesis in \opt{degree}'',
``Author: \opt{candidate}'' and, at the bottom, student ID, advisor(s),
co-advisor(s) and academic year. Chapters open with a large blue number
and bar; starred chapters (\verb|\chapter*{Introduction}|) still enter the
table of contents and the running header. \verb|\blankpage| inserts an
empty page with the raggiera.

\begin{lstlisting}
\documentclass[university=polimi, master, format=classical,
               language=english]{academicthesis}
\usepackage{amssymb, bm, float, colortbl}
\usepackage[colorlinks=true, linkcolor=black, citecolor=black]{hyperref}

\thesissetup{
  title         = {Title},
  candidate     = {Name Surname},
  idnumber      = {000000},
  degree        = {Xxxxxxx Engineering - Ingegneria Xxxxxxx},
  supervisor    = {Prof. Name Surname},
  cosupervisors = {Name Surname, Name Surname},
  academicyear  = {2025-26},
}
\addbibresource{references.bib}

\begin{document}
\frontpage

\frontmatter
\begin{abstractthesis}[keywords = {first, second, third}]
  English abstract.
\end{abstractthesis}
\begin{abstractthesis}[language = italian, keywords = {primo, secondo}]
  Abstract in lingua italiana.
\end{abstractthesis}
\tableofcontents

\mainmatter
\chapter*{Introduction}
\chapter{Chapter one}
...
\printbibliography
\appendix
\chapter{Appendix A}

\backmatter
\listoffigures
\listoftables
\begin{acknowledgements} Thanks to ... \end{acknowledgements}
\end{document}
\end{lstlisting}

\subsection{The \texorpdfstring{\opt{article}}{article} format}

A4, 11pt, 2\,cm margins. There is no separate cover: \cs{frontpage} prints a
header block on page~1 (logo, title, programme, ``author, student ID'',
advisors and academic year beside the \textbf{English abstract}, and the
keywords in a blue box), and the text starts right below. The abstract is
therefore given in \cs{thesissetup}. Use \cs{section}, not \cs{chapter}:
the \pkg{article} class has no chapters (a \cs{chapter} still compiles, as
a section, with a warning).

\begin{lstlisting}
\documentclass[university=polimi, master, format=article,
               language=english]{academicthesis}
\thesissetup{
  title        = {Title},
  candidate    = {Name Surname},
  idnumber     = {000000},
  degree       = {Xxxxxxx Engineering - Ingegneria Xxxxxxx},
  supervisor   = {Prof. Name Surname},
  academicyear = {2025-26},
  abstract     = {The English abstract, printed on page 1.},
  keywords     = {here, the keywords},
}
\begin{document}
\frontpage
\section{Introduction}
...
\appendix
\section{Appendix A}
\clearpage
\begin{abstractthesis}[language = italian, keywords = {qui, le parole chiave}]
  Abstract in lingua italiana.
\end{abstractthesis}
\begin{acknowledgements} ... \end{acknowledgements}
\end{document}
\end{lstlisting}

\subsection{The \texorpdfstring{\opt{summary}}{summary} format (Executive Summary)}

A4, 11pt, two columns. Page~1 opens with a full-width block (logo,
``Executive Summary of the Thesis'', title, programme, author, advisor(s),
co-advisor(s), academic year); the following pages have a blue header with
``Executive summary'' and the author. The student ID is not printed and not
required. As in the template, the summary usually ends with a short
bibliography.

\begin{lstlisting}
\documentclass[university=polimi, master, format=summary,
               language=english]{academicthesis}
\thesissetup{
  title        = {Title of the thesis},
  candidate    = {Name Surname},
  degree       = {Xxxxxxx Engineering - Ingegneria Xxxxxxx},
  supervisor   = {Prof. Name Surname},
  cosupervisor = {Name Surname},
  academicyear = {2025-26},
}
\begin{document}
\frontpage
\section{Introduction}
...
\printbibliography
\end{document}
\end{lstlisting}

\subsection{The \texorpdfstring{\opt{phd}}{phd} format (doctoral dissertation)}

10pt, two-sided, chapters on right-hand pages, 17$\times$24\,cm by default
(set \opt{paper=a4} in \cs{polimithesissetup} for A4), Times for the text and
Helvetica for headings. The cover carries the PoliMi seal, ``Politecnico di
Milano'', the department and doctoral programme, the title, the author on
the right, the supervisor, tutor and chair of the programme on the left,
and the academic year and cycle at the bottom. A \env{defn} (Definition)
environment is provided.

\begin{lstlisting}
\documentclass[university=polimi, phd, language=english]{academicthesis}
\thesissetup{
  title        = {Title},
  candidate    = {Name Surname},
  department   = {Department of Electronics, Information and Bioengineering},
  phdtitle     = {Doctoral Program in Information Technology},
  supervisor   = {Prof. Name Surname},
  tutor        = {Prof. Name Surname},
  coordinator  = {Prof. Name Surname},
  academicyear = {2025--2026},
  cycle        = {XXXVIII},
}
\begin{document}
\frontpage
\frontmatter
\begin{abstractthesis} The abstract. \end{abstractthesis}
\tableofcontents
\mainmatter
\chapter{Introduction}
...
\end{document}
\end{lstlisting}

The \opt{phd} degree selects \opt{format=phd} by itself; \opt{format=phd}
with \opt{bachelor} or \opt{master} is an error.

\subsection{Colours, theorems and other helpers}

\begin{itemize}
  \item The colour \opt{bluepoli} is defined in every format (also as
    \opt{bluePoli} in \opt{article} and \opt{summary}, the name used by the
    official templates): \verb|\rowcolor{bluepoli!40}| (with \pkg{colortbl}).
  \item \env{theorem} and \env{proposition} environments are ready, with
    blue headings; your own \cs{newtheorem}s take the same style.
  \item Figure and table caption labels are blue.
  \item Packages such as \pkg{hyperref}, \pkg{cleveref}, \pkg{amssymb},
    \pkg{bm}, \pkg{float}, \pkg{colortbl}, \pkg{subfig}, \pkg{tikz} and
    \pkg{algorithm} are left to you, as in the official templates.
\end{itemize}

Complete versions are in \file{showcases/example-polimi-classical.tex},
\file{-article.tex}, \file{-summary.tex} and \file{-phd.tex}. The package's
own manual, \file{universities/polimi/academicthesis-polimi-doc.pdf}, gives
further detail on each layout.

% ======================================================================
\section{Common errors and how to fix them}
\label{sec:errors}

\begin{longtable}{@{}L{5.2cm} L{4.2cm} L{4.6cm}@{}}
\toprule
\textbf{Message (abridged)} & \textbf{Cause} & \textbf{Fix} \\
\midrule
\endhead
Missing required thesis metadata: `\textit{key}' & A required field is not in \cs{thesissetup} & Add it; the message shows an example. Tables in Sections~\ref{sec:luiss} and~\ref{sec:polimi}. \\
Missing required thesis metadata: `examdate' & \opt{discussion} without the defence date & Add \opt{examdate}, or remove \opt{discussion}. \\
Unknown format `\textit{x}' for university `polimi' & Typo in \opt{format} & Use \opt{classical}, \opt{article}, \opt{summary} or \opt{phd}. \\
format=phd is the PoliMi doctoral dissertation: use the `phd' degree & \opt{format=phd} with \opt{bachelor} or \opt{master} & Use the \opt{phd} degree. \\
No front page registered for university `\textit{u}', degree `\textit{d}' & A degree the layout does not offer & Check the degree and format options. \\
File `academicthesis.cls' not found & \LaTeX{} cannot see the class files & See Section~\ref{sec:install}: keep the line that inputs \file{academicthesis-repo-paths}, or copy the files next to the thesis. \\
Warning: The base class `article' has no chapters & \cs{chapter} in PoliMi \opt{article} or \opt{summary} & Use \cs{section}. \\
Warning: `french' \dots{} was not loaded by babel & A block with \opt{language=french} & Add \opt{french} to the \opt{other-languages} option. \\
Warning: biblatex is not available & \pkg{biblatex} not installed & Install it, or keep the \pkg{natbib} fallback. \\
\bottomrule
\end{longtable}

Errors from Texifier's built-in engine (for example about \LaTeX3
commands) disappear when the typesetter is set to \TeX{}~Live (pdf\LaTeX).

% ======================================================================
\clearpage
\appendix
\section{Quick reference}
\label{sec:cheatsheet}

\begin{lstlisting}
% ---- LUISS ------------------------------------------------------------
\documentclass[university=luiss, bachelor|master|phd,
               language=english|italian|french,
               no-secondpage, discussion]{academicthesis}
\luissthesissetup{ major = {...}, logo-align = left|center|right,
                   title-align = ..., course-align = ... }
% bachelor/master: title candidate degree chair supervisor(s) academicyear
% phd:             title candidate department phdtitle cycle
%                  supervisor(s) academicyear
% optional:        subtitle idnumber cosupervisor(s) copyyear authoremail
%                  website  (+ examdate examiners with `discussion')

% ---- Politecnico di Milano --------------------------------------------
\documentclass[university=polimi, bachelor|master,
               format=classical|article|summary,
               language=english|italian]{academicthesis}
\documentclass[university=polimi, phd]{academicthesis}       % format=phd
\polimithesissetup{ logo-file = ..., watermark = true|false,
                    paper = 17x24|a4 }                         % paper: phd
% classical: title candidate idnumber degree supervisor(s) academicyear
% article:   + abstract (keywords optional)
% summary:   as classical, without idnumber
% phd:       title candidate department phdtitle supervisor(s)
%            academicyear cycle   (tutor(s), coordinator optional)

% ---- Blocks (every university) ----------------------------------------
\begin{abstractthesis}[language=..., keywords=..., title=...]
  ... \end{abstractthesis}
\begin{originality}[language=..., title=...] ... \end{originality}
\begin{acknowledgements}[language=..., title=...]
  ... \end{acknowledgements}
\dedication[language=...]{...}
\end{lstlisting}

\vfill
{\small\color{docgrey}
\pkg{academicthesis} is released under the LPPL 1.3c or later; copyright
2025--2026 Luca Salvatore. The Politecnico di Milano layouts are based on
the official templates of Politecnico di Milano, whose design, name, logos
and graphic elements belong to Politecnico di Milano; the LUISS name and
logo belong to LUISS Guido Carli. \pkg{academicthesis} is not published by
either university.\par}

\end{document}
